\documentclass{article}
\usepackage[utf8]{inputenc}
\usepackage{amsmath,amssymb}
\usepackage[most]{tcolorbox}
\usepackage{geometry}
\geometry{margin=2.5cm}

\newcommand{\step}[1]{\par\noindent\hangindent=3.5em\hangafter=1%
  \makebox[3.5em][l]{\textbf{#1.}}}
\newcommand{\steptag}[1]{\hfill{\small\textsf{[#1]}}}

\newtcolorbox{proofbox}[1]{
  colback=gray!5, colframe=gray!60,
  fonttitle=\bfseries, title={#1},
  breakable, enhanced,
  left=4pt, right=4pt, top=4pt, bottom=4pt,
  before skip=10pt, after skip=10pt
}

\begin{document}

\begin{proofbox}{Quantitative approximate group laws: there exist universa...}
\step{1} Quantitative approximate group laws: there exist universal C\_grp, C\_pol >= 1, kappa\_pol in (0, 1/2] such that for every finite-dimensional exact-unit epsilon\_r-C*-algebra and delta > 0 satisfying C\_pol*(epsilon\_r + delta) <= kappa\_pol and C\_grp*epsilon\_r < delta - C\_pol*(epsilon\_r*delta + delta\textasciicircum{}2), the inverse u\_delta of the polar map defines C\textasciicircum{}1 maps mu(U, V) = u\_delta(U bold-dot V), sigma(U) = u\_delta(U\textasciicircum{}dagger) on all of calU, with mu(J, U) = mu(U, J) = U, sigma(J) = J, ||mu(U, V) - U bold-dot V|| <= C\_grp*epsilon\_r, ||sigma(U) - U\textasciicircum{}dagger|| <= C\_grp*epsilon\_r, ||mu(mu(U, V), W) - mu(U, mu(V, W))|| <= C\_grp*epsilon\_r, ||mu(sigma(U), U) - J|| <= C\_grp*epsilon\_r, and ||mu(U, sigma(U)) - J|| <= C\_grp*epsilon\_r. \steptag{validated} \\[6pt]
\step{1.1} Choose one universal constant triple and one common fixed-delta polar datum as follows: if (G\_d,P\_d,k\_d), (G\_c,P\_c,k\_c), and (P\_r,k\_r) are witnesses furnished respectively by lem-stage1-group-domain-membership, lem-stage1-group-closeness, and lem-stage1-polar-retraction, set C\_pol=max(P\_d,P\_c,P\_r), kappa\_pol=min(k\_d,k\_c,k\_r,1/16), and C\_grp=max(G\_d,8*G\_c,8). For every algebra and delta satisfying the root guards, all three upstream results apply, their polar inverse data can be identified, every U bold-dot V and U\textasciicircum{}dagger lies in the resulting S\_delta, and the common inverse obeys the two bounds with constant G\_c. \steptag{validated} \\[6pt]
\step{1.1.1} Let (G\_d,P\_d,k\_d), (G\_c,P\_c,k\_c), and (P\_r,k\_r) be fixed universal witnesses for lem-stage1-group-domain-membership, lem-stage1-group-closeness, and lem-stage1-polar-retraction. Define P=max(P\_d,P\_c,P\_r), k=min(k\_d,k\_c,k\_r,1/16), and G=max(G\_d,8*G\_c,8). Then P,G>=1 and 0<k<=1/2. If P*(epsilon\_r+delta)<=k and G*epsilon\_r<delta-P*(epsilon\_r*delta+delta\textasciicircum{}2), then P\_i*(epsilon\_r+delta)<=k\_i for i=d,c,r, while G\_i*epsilon\_r<=G*epsilon\_r<delta-P*(epsilon\_r*delta+delta\textasciicircum{}2)<=delta-P\_i*(epsilon\_r*delta+delta\textasciicircum{}2) for i=d,c. Hence all hypotheses of those three exact named externals hold; also epsilon\_r<=1/16. \steptag{validated} \\[4pt]
\step{1.1.2} For the fixed algebra and delta, every polar datum produced above uses the same map Pi\_delta:(U,H)$|-\rangle$U bold-dot H on the same full domain calU x B\_delta\textasciicircum{}\{calH\}(J), so every associated S\_delta is its set-theoretic image and these image sets coincide. The external lem-stage1-polar-coherence-naturality identifies their inverse components on that common image. Therefore the membership conclusions of lem-stage1-group-domain-membership, the G\_c*epsilon\_r estimates of lem-stage1-group-closeness, and the C\textasciicircum{}1 inverse supplied by lem-stage1-polar-retraction all concern one common u\_delta:S\_delta->calU; in particular U bold-dot V and U\textasciicircum{}dagger belong to S\_delta for all U,V in calU. \steptag{validated} \\[4pt]
\step{1.2} By 1.1, U bold-dot V and U\textasciicircum{}dagger lie in the common open S\_delta and u\_delta:S\_delta->calU is C\textasciicircum{}1 by lem-stage1-polar-retraction. The multiplication is continuous bilinear by def-epsilon-cstar-algebra and hence C\textasciicircum{}1 in finite dimension; dagger is conjugate-linear, hence real-linear and C\textasciicircum{}1. Therefore mu(U,V)=u\_delta(U bold-dot V) and sigma(U)=u\_delta(U\textasciicircum{}dagger) are globally defined C\textasciicircum{}1 maps into calU. Also J lies in calU, J bold-dot U=U bold-dot J=U, and J\textasciicircum{}dagger=J by the exact-unit clauses. Since lem-stage1-polar-retraction gives u\_delta(T)=T for every T in calU, mu(J,U)=mu(U,J)=U and sigma(J)=J. \steptag{validated} \\[4pt]
\step{1.3} For the common u\_delta fixed in 1.1, lem-stage1-group-closeness gives ||mu(U,V)-U bold-dot V||<=G\_c*epsilon\_r and ||sigma(U)-U\textasciicircum{}dagger||<=G\_c*epsilon\_r. Since C\_grp=G>=8*G\_c, each is at most C\_grp*epsilon\_r, including epsilon\_r=0. \steptag{validated} \\[4pt]
\step{1.4} The guard in 1.1 gives epsilon\_r<=1/16. For T in calU, def-approximate-unitary-space gives T\textasciicircum{}dagger bold-dot T=J, while the exact-unit clause of def-epsilon-cstar-algebra gives ||J||=1. Its C*-inequality yields (1-epsilon\_r)||T||\textasciicircum{}2<=||T\textasciicircum{}dagger bold-dot T||=1. Hence ||T||<=sqrt(16/15)<4/3 (because 16/15<16/9). This also bounds ||T\textasciicircum{}dagger||=||T||. \steptag{validated} \\[4pt]
\step{1.5} Put A=mu(U,V) and B=mu(V,W), which lie in calU by 1.2. Insert A bold-dot W, (U bold-dot V) bold-dot W, U bold-dot (V bold-dot W), and U bold-dot B between the two iterated products. The two outer differences are each <=G\_c*epsilon\_r by lem-stage1-group-closeness; bilinearity and the product-norm axiom bound each replacement difference by (1+epsilon\_r)*(4/3)*G\_c*epsilon\_r<=(17/12)*G\_c*epsilon\_r using 1.4; and the sole reassociation difference is <=epsilon\_r*(4/3)\textasciicircum{}3=(64/27)*epsilon\_r by the associator axiom. Thus the total is <=(2+17/6)*G\_c*epsilon\_r+(64/27)*epsilon\_r<=(389/54)*G\_c*epsilon\_r<=8*G\_c*epsilon\_r<=C\_grp*epsilon\_r, where G\_c>=1. All inequalities remain non-strict at epsilon\_r=0. \steptag{validated} \\[4pt]
\step{1.6} Fix U in calU and choose a right inverse R, which exists by def-approximate-unitary-space. Then ||U bold-dot U\textasciicircum{}dagger-J||<=4*epsilon\_r. This estimate is derived locally below from the root smallness guard and the defining epsilon\_r-C*-algebra axioms, without using sibling node 1.4. \steptag{validated} \\[6pt]
\step{1.6.1} The existential constants in the root proof may be, and are, chosen with kappa\_pol<=1/16: shrinking a positive admissible kappa\_pol to min(kappa\_pol,1/16) preserves all upstream hypotheses and conclusions. Hence the root guard C\_pol*(epsilon\_r+delta)<=kappa\_pol, together with C\_pol>=1 and delta>0, gives epsilon\_r<=epsilon\_r+delta<=kappa\_pol/C\_pol<=1/16. For U in calU, def-approximate-unitary-space gives U\textasciicircum{}dagger bold-dot U=J, while the exact-unit clause of def-epsilon-cstar-algebra gives ||J||=1. The C*-inequality therefore yields (1-epsilon\_r)*||U||\textasciicircum{}2<=||U\textasciicircum{}dagger bold-dot U||=1, so ||U||\textasciicircum{}2<=16/15<16/9 and ||U||=||U\textasciicircum{}dagger||<=4/3. \steptag{validated} \\[4pt]
\step{1.6.2} Choose the right inverse R supplied by def-approximate-unitary-space, so U bold-dot R=J. Exact-unit identities and U\textasciicircum{}dagger bold-dot U=J give R-U\textasciicircum{}dagger=(U\textasciicircum{}dagger bold-dot U) bold-dot R-U\textasciicircum{}dagger bold-dot (U bold-dot R). The associator axiom and involutive isometry give ||R-U\textasciicircum{}dagger||<=epsilon\_r*||U\textasciicircum{}dagger||*||U||*||R||=a*||R||, where a:=epsilon\_r*||U||\textasciicircum{}2. By the preceding local estimate, a<=(16/9)*epsilon\_r<=1/9. Thus ||R||<=||U\textasciicircum{}dagger||+||R-U\textasciicircum{}dagger||<=4/3+a*||R||, whence (1-a)||R||<=4/3 and ||R||<=(4/3)/(8/9)=3/2. \steptag{validated} \\[4pt]
\step{1.6.3} By bilinearity and U bold-dot R=J, U bold-dot U\textasciicircum{}dagger-J=U bold-dot (U\textasciicircum{}dagger-R). The product-norm axiom, followed by the two preceding local estimates, gives ||U bold-dot U\textasciicircum{}dagger-J||<=(1+epsilon\_r)*||U||*||U\textasciicircum{}dagger-R||<=(1+epsilon\_r)*||U||*epsilon\_r*||U||\textasciicircum{}2*||R||<=(17/16)*(4/3)*(16/9)*(3/2)*epsilon\_r=(34/9)*epsilon\_r<=4*epsilon\_r. This uses only bilinearity, the product-norm axiom, and the single associator estimate in the preceding child. \steptag{validated} \\[4pt]
\step{1.7} Since sigma(U) lies in calU, lem-stage1-group-closeness applies to (sigma(U),U) and gives ||mu(sigma(U),U)-sigma(U) bold-dot U||<=G\_c*epsilon\_r. By the product-norm axiom, 1.4, and ||sigma(U)-U\textasciicircum{}dagger||<=G\_c*epsilon\_r, the next difference is ||sigma(U) bold-dot U-U\textasciicircum{}dagger bold-dot U||<=(17/12)*G\_c*epsilon\_r. Finally U\textasciicircum{}dagger bold-dot U=J by def-approximate-unitary-space. Hence ||mu(sigma(U),U)-J||<=(29/12)*G\_c*epsilon\_r<=3*G\_c*epsilon\_r<=C\_grp*epsilon\_r, also at epsilon\_r=0. \steptag{validated} \\[4pt]
\step{1.8} Since sigma(U) lies in calU, lem-stage1-group-closeness gives ||mu(U,sigma(U))-U bold-dot sigma(U)||<=G\_c*epsilon\_r. The product-norm axiom, 1.4, and ||sigma(U)-U\textasciicircum{}dagger||<=G\_c*epsilon\_r give ||U bold-dot sigma(U)-U bold-dot U\textasciicircum{}dagger||<=(17/12)*G\_c*epsilon\_r. Node 1.6 gives ||U bold-dot U\textasciicircum{}dagger-J||<=4*epsilon\_r. Therefore ||mu(U,sigma(U))-J||<=(29/12)*G\_c*epsilon\_r+4*epsilon\_r<=(77/12)*G\_c*epsilon\_r<=8*G\_c*epsilon\_r<=C\_grp*epsilon\_r because G\_c>=1; this is endpoint-safe at epsilon\_r=0. \steptag{validated} \\[4pt]
\end{proofbox}

\end{document}
